\subsection{Channel Models for Simulation}
\label{sec:modelo-canal}
Evaluating a massive MIMO system requires a model that represents the channel while capturing the key characteristics of an array with a large number of antennas. Such a model must account for the antenna array geometry, the spatial correlation between the channel responses of different antennas, and the physical characteristics of the environment—such as the locations and orientations of base stations (BSs) and user equipment (UEs) \cite{bjornson-mimo}.

\subsubsection{Types of Channel Models}
There are two main types of channel models: deterministic and stochastic. A deterministic model always yields the same result under identical conditions, as it relies on a specific environment with fixed parameters—such as the locations of transmitters, receivers, signal scatterers, and reflectors. Examples include ray-tracing and empirical, field-based channel measurements. A drawback of this approach is that reproducing results often requires access to specific channel measurement data or 3D building models, which are not always available \cite{bjornson-mimo}.

Stochastic models are independent of any particular environment and generalize channel behavior based on selected statistical properties. They are divided into three categories \cite{bjornson-mimo}:

\begin{itemize}
\item Correlation-based: Channel responses follow a zero-mean Gaussian distribution, defined via correlation matrices. One example is the correlated Rayleigh fading model, which assumes frequency-flat fading—meaning all multipath components arrive at the antenna with the same delay. This model is sufficient for analysis within a coherence block (Section~\ref{subsec:bloco-coerencia}). 
\item Parametric: These models define statistical distributions for the number of scatterers, delays, powers, and angles of arrival/departure (AoA/AoD). They do not account for the environment's geometry and, therefore, do not cover the analysis of the channel's temporal evolution resulting from UE movement. 
\item Geometry-based: These models strike a balance between the characteristics of stochastic and deterministic models. They define a spatial distribution for the physical locations of signal scatterers surrounding transmitters and receivers and, based on this, model the channel's propagation paths in a quasi-deterministic manner. A quasi-deterministic model defines a process governed by fixed laws, yet one that incorporates a degree of randomness through noise and/or variability.
\end{itemize}

% FONTE: Björnson, Hoydis e Sanguinetti 2017, Seção 7.3, pp. 329-331 (Fig. 7.10).
